\section{Experimental Setup}
\label{sec:setup}

\subsection{Dataset and Evaluation Splits}
\label{sec:setup-data}

We evaluate on \tool{}~\cite{ma2026autodojo}, a trajectory-level benchmark of tool-using agent executions containing 2{,}218 trajectories spanning 2{,}000 tasks in five domains, of which 697 are compromised and 1{,}521 are clean. Each trajectory is converted into a typed directed acyclic graph with four node types (\textsc{Problem}, \textsc{Thinking}, \textsc{Action}, \textsc{State}) and five relation types (\texttt{depends\_on}, \texttt{computes\_from}, \texttt{reflects}, \texttt{sequential}, \texttt{branches\_from}), following the construction in Section~\ref{sec:method-graphs}. The benchmark provides an official split into 1{,}253 training, 267 validation, and 269 test trajectories plus a pool of 429 compromised trajectories that are excluded from the supervised splits; we deliberately reserve this unused pool for the defense experiments in Section~\ref{sec:setup-defense}, so that no case whose experience is synthesized was ever seen by the detector during training, validation, or threshold calibration.

For detector training we follow a stratified 7:1.5:1.5 protocol in which every split contains 15\% abnormal graphs (split seed 2026). This differs from the official split because the official test set is drawn from the same task distribution as training, whereas our protocol requires disjoint, identically distributed partitions for multi-seed variance estimation. All label access is confined to the training and validation splits; the unused pool is queried only after detector training is complete and is used exclusively through the defense pipeline, in which ground-truth labels never enter the LLM context.

\subsection{Detector Configuration}
\label{sec:setup-detector}

We train the \method{} detector as an edge-attention R-GCN~\cite{schlichtkrull2017modeling,velikovi2017graph} over frozen sentence embeddings~\cite{reimers2019sentencebert} (all-MiniLM-L6-v2, 384~d; a parallel variant encodes directed edges as concatenations of source and target text). The configuration matches the strongest hyperparameter setting identified on this dataset in prior runs: hidden dimension 64, two R-GCN layers with four attention heads, relation and direction embeddings of 16 and 8 dimensions respectively, mean--max readout pooling, cross-entropy classification loss, dropout $0.02$, learning rate $7\times10^{-4}$, weight decay $10^{-4}$, batch size 8, up to 120 epochs with early stopping (patience 35). Injection-pattern keywords are masked from node text during both training and inference (\texttt{mask-leakage-words}), so the detector cannot rely on lexical shortcuts such as the verbatim attack string. The operating threshold is selected on the validation split by maximizing abnormal-class F$_1$; we report detection quality as ROC-AUC, PR-AUC, accuracy, recall, and abnormal F$_1$ at that threshold. We train five seeds (42--46) and report mean~$\pm$~standard deviation; all runs execute on a single NVIDIA RTX 4090 (24~GB) with PyTorch~2.7.1 and PyTorch-Geometric~2.6.1.

The node risk scorer is a two-layer MLP over $[\mathbf{h}_i \,\|\, \mathbf{h}_G \,\|\, \mathbf{t}_i]$ trained against the recording-grounded supervision of Section~\ref{sec:method-localization} on the seed-42 detector's training-split predictions, with a validation-set optimal threshold chosen by F$_1$ on compromised nodes. At inference we score every node of a flagged graph, take the top-3 nodes by risk, and extract their BFS neighborhoods as the risk subgraph that conditions experience synthesis.

\subsection{Defense Evaluation Protocol}
\label{sec:setup-defense}

The main evaluation is a paired, within-case comparison on 100 attack trajectories sampled from the unused pool (sample seed 20261006; sampling is uniform over the pool without stratification). Each trajectory is replayed twice by the same agent model under identical decoding configuration: once as a \emph{baseline} run with no injected experience, and once as an \emph{experience} run preceded by the safety memory synthesized for that trajectory by the detection-triggered pipeline of Section~\ref{sec:method-experience}. Both runs are judged for attack success by an independent evaluator model that is different from the agent model, closing the self-preference channel that recent analyses of LLM-as-judge pipelines identify as a first-order threat~\cite{zhou2026style,xu2026inside,weng2026beyond}. Judge--agent agreement is additionally spot-checked against the recorded compromised execution that each benchmark trajectory carries, and verdict variance under fixed decoding is monitored in light of known within-judge noise floors~\cite{ayyagari2026pinned}.

Model roles are fixed across the study and published in full. The \emph{gate} controller that adjudicates every replayed step (the experience-conditioned replay agent of Section~\ref{sec:method-replay}) runs GPT-5.5-Fast; the \emph{experience synthesis judge} runs DeepSeek-V4-Pro for the main strict-gate panel---a deliberate substitution made after pilot runs showed that the originally planned Flash variant's unbounded reasoning loops truncated its JSON verdicts on 41\% of synthesis calls even at a 16{,}384-token completion budget, whereas the Pro variant completed every pilot call---and GPT-5.5 for the recall-gate delta and the ablation experience variants synthesized after the DeepSeek endpoint's balance was exhausted mid-study (357 of the 413 strict-panel judgments, 78 of the 79 recall-delta judgments, and all ablation variants come from their respective judges, with the judge identity stored per judgment row); and the \emph{attack-success judge} runs GPT-5.5 through a separate endpoint, ensuring that neither the agent nor the synthesizer grades its own output. The primary endpoint is the attack success rate (ASR) difference between paired conditions, assessed with McNemar's exact test on the discordant pairs and the Wilcoxon signed-rank test as a robustness check; we also report per-domain breakdowns and the fraction of trajectories whose replay terminates in a safe action class. Detection quality (Section~\ref{sec:setup-detector}) is reported as a secondary endpoint on the held-out test splits of all five seeds.

\subsection{Ablations and Analyses}
\label{sec:setup-ablation}

Four ablations isolate the pipeline's contribution structure: (i)~\emph{raw-trajectory experience}, in which the synthesis judge receives the full trajectory text instead of the localized risk subgraph, testing whether graph-based localization improves experience specificity; (ii)~\emph{no-localization}, in which the entire flagged graph is rendered for the judge; (iii)~\emph{generic safety prompt}, a non-adaptive control in which a fixed safety instruction of matched length replaces the injected experience; and (iv)~\emph{coverage}, in which the detector's operating threshold is moved to its validation recall-weighted F$_2$ optimum so that far more cases trigger synthesis---the recall-oriented gate of Section~\ref{sec:results-ablation}---while fully unflagged cases reduce to the baseline by construction (no flag, no experience). We additionally analyze cross-case transfer (experiences synthesized from one attack applied to held-out attacks of the same family), report the sensitivity of the defense to the detector's threshold-selection rule by contrasting the validation-F$_1$ and validation-F$_2$ operating points of the same scorer, and examine the risk subgraphs and experience fields of representative success and failure cases.
